\section{A Gaussian collision test through two theta primitives}
\label{sec:primitive}

The complete source permits an exact transfer of a quadratic collision
observation to undifferentiated theta primitives. We incorporate the
two-source calculation of \cite{TNDPrimitiveLattice}, using the radial
identity and boundary accounting of \cite{RadialThetaPrimitive,RadialThetaIBP}.
The integer-index polynomial factors disappear from the primitive product;
derivatives act on an explicit test, and a boundary curvature term remains.
These identities retain proof-author checking. Their signed estimating
consequence is still open.

\subsection{The Gaussian observation and its source kernel}
For real $\lambda,x$ and $v>0$, put
\begin{align}
 w_v(s)&=(2\pi v)^{-1/2}\ee^{-s^2/(2v)},\qquad
 F_\lambda(u)=\ee^{\lambda u^2}\Phi(u),\label{eq:primitiveweight}\\
 A_0(\lambda,x,v)&=\int_\R w_v(k-x)
       [H_\lambda'(k)^2-H_\lambda(k)H_\lambda''(k)]\dd k.
 \label{eq:A0definition}
\end{align}
Here $x$ is a scalar observation center, distinct from the position operator
$X$, and primes on $H$ denote readout derivatives. The identities hold at
every finite real $\lambda$; this does not extend the physical apparatus
beyond its specified time slab. On that slab, with $\alpha=\gamma t>0$
and $h_*=H_{1/2}(0)$, \eqref{eq:dj} gives
\begin{equation}\label{eq:A0count}
 A_0=h_*^2\int_\R w_v(k-x)
       [(\partial_kd_0)^2+\alpha^{-1}d_0d_1]\dd k.
\end{equation}
This is a product of analytical traces, not a joint probability for two
different count outcomes in one trial.

\begin{proposition}[Complete source pairing]\label{prop:A0source}
Define $g_{x,v}(s)=\ee^{-vs^2/2}\cos(xs)$ and
\begin{align}
 K_{x,v}(p,q)&=\frac14[(p-q)^2g_{x,v}(p+q)
                         +(p+q)^2g_{x,v}(p-q)],\label{eq:A0K}\\
 T_{\lambda,x,v}(p,q)&=\ee^{\lambda(p^2+q^2)}K_{x,v}(p,q).
 \label{eq:A0T}
\end{align}
Then
\begin{align}
 A_0&=\frac18\iint_{\R^2}(p-q)^2F_\lambda(p)F_\lambda(q)
       g_{x,v}(p+q)\dd p\dd q,\label{eq:A0fullsource}\\
 A_0&=\int_0^\infty\!\int_0^\infty
       \Phi(p)\Phi(q)T_{\lambda,x,v}(p,q)\dd p\dd q.
 \label{eq:A0folded}
\end{align}
The test $T$ is symmetric in $p,q$ and even in each variable separately.
\end{proposition}
\begin{proof}
Evenness gives $H_\lambda(k)=\tfrac12\int_\R F_\lambda(u)\ee^{\ii ku}\dd u$.
In the double integral for $H'^2-HH''$, symmetrizing $-pq+q^2$ gives
$(p-q)^2/2$. The two field factors $1/2$ produce the coefficient $1/8$.
Gaussian integration uses the exact identity
$\int_\R w_v(k-x)\ee^{\ii ks}\dd k=\ee^{\ii xs-vs^2/2}$.
Taking real parts proves \eqref{eq:A0fullsource}; reflection of the four
quadrants proves \eqref{eq:A0folded}. Complete source moments justify all
interchanges. No fluctuating variable has been replaced by its mean.
\end{proof}

\subsection{The radial primitive and its unit-sphere trace}
For $u\ge0$, set
\begin{equation}\label{eq:radialprimitive}
 r=\ee^{2u},\qquad A(r)=\sum_{n\ge1}\ee^{-\pi n^2r^2},\qquad
 \psi(u)=\ee^u A(\ee^{2u}),\qquad D=\partial_u^2-1.
\end{equation}
The radial Laplacian in three dimensions satisfies
\begin{equation}\label{eq:radialsource}
 \Delta_3=\partial_r^2+2r^{-1}\partial_r,\qquad
 \Phi(u)=\tfrac12r^{5/2}\Delta_3A(r),\qquad D\psi=8\Phi.
\end{equation}
Indeed $\Delta_3\ee^{-cr^2}=(4c^2r^2-6c)\ee^{-cr^2}$ and the
change $r=\ee^{2u}$ gives $D\psi=4r^{5/2}\Delta_3A$.
Spatial dimension three does not change the rank-one arithmetic index.

Write $\Theta(r)=1+2A(r)$. The modular identity
$\Theta(r)=r^{-1}\Theta(1/r)$, the zero-argument specialization of
\cite[20.7.32]{DLMFThetaModular}, implies
\begin{equation}\label{eq:primitiveendpoint}
 A'(1)=-\tfrac12A(1)-\tfrac14,\qquad
 \psi(0)=A(1),\qquad \psi'(0)=-\tfrac12.
\end{equation}
For a twice differentiable test $f$ whose derivatives are dominated by
the complete theta tail, integration by parts gives
\begin{equation}\label{eq:primitiveGreen}
 \int_0^\infty\Phi f\dd u
 =\frac{f(0)}{16}+\frac{A(1)f'(0)}8
   +\frac18\int_0^\infty\psi(f''-f)\dd u.
\end{equation}
The boundary is $[\psi'f-\psi f']_0^\infty$. Equivalently, with
$h(r)=r^{-1/2}f(\log r/2)$, the right side is
\begin{equation}\label{eq:primitiveRadialGreen}
 \frac14\int_1^\infty A(r)\Delta_3h(r)r^2\dd r
 +\frac14\{-A'(1)h(1)+A(1)h'(1)\}.
\end{equation}
Thus the endpoint is the unit-sphere contribution of the three-dimensional
Green identity. The positive-time factor remains in the test, leaving the
rapidly decreasing primitive unchanged.

\subsection{Transferring both source factors}
\begin{theorem}[Two primitives and the complete boundary]\label{thm:twoPrimitive}
Let $D_p=\partial_p^2-1$, $D_q=\partial_q^2-1$, and let $T$ be
\eqref{eq:A0T}. Set
\begin{equation}\label{eq:primitiveface}
 f(u)=T(0,u)=\tfrac12u^2\ee^{(\lambda-v/2)u^2}\cos(xu).
\end{equation}
All integrals below converge absolutely, and
\begin{align}
 A_0={}&\frac1{64}\left[
  \iint_{[0,\infty)^2}\psi(p)\psi(q)D_pD_qT(p,q)\dd p\dd q
  +\int_0^\infty\psi(u)(f''-f)(u)\dd u\right],
  \label{eq:twoPrimitiveFace}\\
 ={}&\frac1{64}\iint_{[0,\infty)^2}
       \psi(p)\psi(q)D_pD_qT(p,q)\dd p\dd q
       -\frac1{16}H_{\lambda-v/2}''(x).
  \label{eq:twoPrimitiveCurvature}
\end{align}
\end{theorem}
\begin{proof}
Apply \eqref{eq:primitiveGreen} in each variable of \eqref{eq:A0folded}.
Separate evenness of $T$, preserved by the other variable's $D$ operator,
makes every first-derivative endpoint trace zero. The remaining terms are
\begin{align*}
 64A_0={}&\iint\psi(p)\psi(q)D_pD_qT(p,q)\dd p\dd q\\
 &+\tfrac12\int\psi(q)D_qT(0,q)\dd q
  +\tfrac12\int\psi(p)D_pT(p,0)\dd p+\tfrac14T(0,0).
\end{align*}
The corner vanishes because $T(0,0)=0$, and symmetry combines the two
faces into the last integral in \eqref{eq:twoPrimitiveFace}. Since
$f(0)=f'(0)=0$, another use of \eqref{eq:primitiveGreen} gives
\[
 \int\psi(f''-f)\dd u=8\int\Phi f\dd u
 =4\int_0^\infty u^2\ee^{(\lambda-v/2)u^2}\Phi(u)\cos(xu)\dd u
 =-4H_{\lambda-v/2}''(x),
\]
proving \eqref{eq:twoPrimitiveCurvature}.

For convergence, $n^2-1\ge3(n-1)$ gives, for $r\ge1$ and $u\ge0$,
\begin{equation}\label{eq:primitiveDecay}
 A(r)\le\frac{\ee^{-\pi r^2}}{1-\ee^{-3\pi r^2}},\qquad
 \psi(u)\le\frac{\ee^{u-\pi\ee^{4u}}}{1-\ee^{-3\pi}}.
\end{equation}
Derivatives of $\psi$ have the same double-exponential decay with fixed
additional factors. On compact real parameter sets with $v\ge0$, all
derivatives of $T$ used here are bounded by
$C(1+p+q)^6\ee^{L(p^2+q^2)}$, where $L=\max(0,\sup\lambda)$.
Equation~\eqref{eq:primitiveDecay} pays every derivative transfer,
vanishing far-end boundary and interchange, and permits $v\downarrow0$.
\end{proof}

For $a=p+q$, $b=p-q$, put $h(s)=s^2\ee^{\lambda s^2/2}$ and
$j(s)=\ee^{(\lambda-v)s^2/2}\cos(xs)$. The transferred test is explicit:
\begin{align}
 T&=\tfrac14\{h(b)j(a)+h(a)j(b)\},\label{eq:primitiveSeparatedTest}\\
 D_pD_q&=\partial_a^4+\partial_b^4-2\partial_a^2\partial_b^2
          -2\partial_a^2-2\partial_b^2+1.
 \label{eq:primitiveTransferredOperator}
\end{align}
The quadrant becomes $a\ge|b|$, with $\dd p\dd q=\dd a\dd b/2$;
the separated test does not give a product integration domain.

\subsection{The exponential lattice and complete truncation payment}
The primitive product is exactly
\begin{equation}\label{eq:primitiveLattice}
 \psi(p)\psi(q)=\ee^{p+q}\sum_{m,n\ge1}
      \exp\{-\pi(m^2\ee^{4p}+n^2\ee^{4q})\}.
\end{equation}
All square-index polynomial prefactors have moved into
\eqref{eq:primitiveTransferredOperator}. In radial coordinates,
\begin{equation}\label{eq:primitiveThetaProduct}
 A(r)A(s)=\sum_{m,n\ge1}\ee^{-\pi(m^2r^2+n^2s^2)}
 =\tfrac14\{\Theta(r)\Theta(s)-\Theta(r)-\Theta(s)+1\}.
\end{equation}
This is the positive-quadrant part of a rectangular rank-two Gaussian
lattice. The two spatial operators can be placed on $\R^3\times\R^3$
without replacing the arithmetic by a rank-six lattice.

Set $\widehat T(r,s)=T(\log r/2,\log s/2)$,
$\widehat f(r)=f(\log r/2)$ and $\mathscr D_r=(2r\partial_r)^2-1$.
The radial version of \eqref{eq:twoPrimitiveFace} is
\begin{align}
 A_0={}&\frac1{256}\int_1^\infty\!\int_1^\infty
 \frac{A(r)A(s)}{\sqrt{rs}}\mathscr D_r\mathscr D_s\widehat T(r,s)
       \dd r\dd s\notag\\
 &+\frac1{128}\int_1^\infty\frac{A(r)}{\sqrt r}
                  \mathscr D_r\widehat f(r)\dd r.
 \label{eq:primitiveRadialPair}
\end{align}
Absolute convergence permits summation under this integral. For integers $N\ge0$,
\begin{equation}\label{eq:primitiveTail}
 R_N(r):=\sum_{n>N}\ee^{-\pi n^2r^2}
 \le\frac{\ee^{-\pi(N+1)^2r^2}}
           {1-\ee^{-\pi(2N+3)r^2}}\qquad(r\ge1).
\end{equation}
Compare successive square gaps with $2N+3$ to prove this bound. With
$A_N=A-R_N$, the omitted product is at most
$R_N(r)A(s)+A(r)R_N(s)$. Its absolute integral against the bulk test in
\eqref{eq:primitiveRadialPair}, divided by 256, plus the face integral
with $R_N$, divided by 128, pays the complete truncation. This approximates
the transferred canonical formula with its exact endpoint coefficient.
A finite $A_N$ does not itself inherit the modular endpoint identity.
No numerical cutoff or margin is selected.

The transformation
$\Theta(r)\Theta(s)=(rs)^{-1}\Theta(1/r)\Theta(1/s)$ is available.
The four terms in braces in \eqref{eq:primitiveThetaProduct} must remain
combined: their separate large-radius integrals against the growing test
need not converge. The time weight and the logarithmic cosine also remain.

\subsection{Normalization, cancellation and the open estimate}
Take $\lambda=x=0$ and $v\downarrow0$, and write
$J_j=\int_0^\infty u^j\psi(u)\dd u$. The transfer gives
\begin{equation}\label{eq:primitiveNormalization}
 A_0=\frac{J_0J_2-2J_0^2+J_0-J_2/2}{64}
 =\frac{1/2-J_0}{8}\,\frac{2J_0-J_2}{8}.
\end{equation}
The last factors are $\int\Phi$ and $\int u^2\Phi$ from
\eqref{eq:primitiveGreen}, as required by $H'(0)^2-H(0)H''(0)$.
Omitting the face term fails this check.

The boundary is not generally small. If
$H_{\lambda_*}(k_*)=H_{\lambda_*}'(k_*)=0$ and
$a=H_{\lambda_*}''(k_*)\ne0$, then along
$(\lambda,x,v)=(\lambda_*-\varepsilon,k_*,\varepsilon)$ it tends to
$-a/16$, while $A_0$ tends to zero. On the initial event box in
\cite{TNDB03Gaussian}, the complete-source diagnostic gives
$A_0=-a^2\varepsilon/2+O(\varepsilon^{3/2})$, with an explicit
curvature-dependent error there; no uniform curvature floor is assumed.
Bulk and boundary must be controlled jointly to preserve this scale.

Positive primitives alone do not separate theta from the collision
controls. For any suitable positive half-line source $q$,
\[
 \psi_q(u)=8\int_u^\infty\sinh(s-u)q(s)\dd s>0,
 \qquad \psi_q''-\psi_q=8q;
\]
a positive rescaling can arrange $\psi_q'(0)=-1/2$ without removing a
double zero \cite{RadialThetaIBP}. Reflected and positive-primitive
controls therefore test the generic transfer but lack the complete
unit-coefficient primitive in \eqref{eq:primitiveThetaProduct}.
Weighted-square and finite-prime controls can retain selected arithmetic
while failing complete coefficients or full reflection. Euler-factor
controls change the literal arithmetic and conductor. Fixed-amplitude,
variable-kernel controls instead fail the cosine/Newman test
\eqref{eq:A0T}. These are failed exact premises, not proof that a future
quantitative inequality fails for each control.

The established simplification removes square-index polynomial weights
while retaining every derivative and boundary contribution. A
source-specific signed estimate of their sum remains open. Direct scalar
analysis may use the same identity; the representation alone establishes
no comparison advantage or collision restriction.
